\begin{thebibliography}{99}
\bibitem{mustovicary} B. Musto y J. Vicary, ``Quantum Latin Squares and Unitary Error Bases'', \emph{Quantum Information and Computation} \textbf{16}, 1318--1332 (2016). \href{https://doi.org/10.26421/QIC16.15-16-4}{doi:10.26421/QIC16.15-16-4}.

\bibitem{delascuevas2020} G. De las Cuevas, T. Drescher y T. Netzer, ``Quantum Magic Squares: Dilations and Their Limitations'', \emph{Journal of Mathematical Physics} \textbf{61}, 111704 (2020). \href{https://doi.org/10.1063/5.0022344}{doi:10.1063/5.0022344}.

\bibitem{delascuevas2023} G. De las Cuevas, T. Netzer e I. Valentiner-Branth, ``Magic Squares: Latin, Semiclassical, and Quantum'', \emph{Journal of Mathematical Physics} \textbf{64}, 022201 (2023). \href{https://doi.org/10.1063/5.0127393}{doi:10.1063/5.0127393}.

\bibitem{codata2022} P. J. Mohr, D. B. Newell, B. N. Taylor y E. Tiesinga, ``CODATA Recommended Values of the Fundamental Physical Constants: 2022'', \emph{Reviews of Modern Physics} \textbf{97}, 025002 (2025). \href{https://doi.org/10.1103/RevModPhys.97.025002}{doi:10.1103/RevModPhys.97.025002}. \href{https://physics.nist.gov/cuu/Constants/Table/allascii.txt}{Tabla de constantes}.

\bibitem{flm} I. B. Frenkel, J. Lepowsky y A. Meurman, \emph{Vertex Operator Algebras and the Monster}, Pure and Applied Mathematics, vol. 134, Academic Press, 1988.

\bibitem{borcherds} R. E. Borcherds, ``Monstrous Moonshine and Monstrous Lie Superalgebras'', \emph{Inventiones Mathematicae} \textbf{109}, 405--444 (1992). \href{https://doi.org/10.1007/BF01232032}{doi:10.1007/BF01232032}.

\bibitem{conwaysloane} J. H. Conway y N. J. A. Sloane, \emph{Sphere Packings, Lattices and Groups}, 3.\textsuperscript{a} ed., Grundlehren der mathematischen Wissenschaften 290, Springer, 1999. \href{https://doi.org/10.1007/978-1-4757-6568-7}{doi:10.1007/978-1-4757-6568-7}.

\bibitem{bipm2018} Conférence générale des poids et mesures, \emph{Resolution 1 of the 26th CGPM: On the revision of the International System of Units}, 2018. \href{https://www.bipm.org/en/committees/cg/cgpm/26-2018/resolution-1}{Resolución oficial}. \href{https://doi.org/10.59161/CGPM2018RES1E}{doi:10.59161/CGPM2018RES1E}.

\bibitem{Bose1924} Satyendra~Nath Bose.
\newblock Plancks gesetz und lichtquantenhypothese.
\newblock {\em Zeitschrift für Physik}, 26:178--181, 1924.

\bibitem{Dirac1926} Paul A.~M. Dirac.
\newblock On the theory of quantum mechanics.
\newblock {\em Proceedings of the Royal Society A}, 112:661--677, 1926.

\bibitem{Einstein1924} Albert Einstein.
\newblock Quantentheorie des einatomigen idealen gases.
\newblock {\em Sitzungsberichte der Preussischen Akademie der Wissenschaften},
  pages 261--267, 1924.

\bibitem{Fermi1926} Enrico Fermi.
\newblock Sulla quantizzazione del gas perfetto monoatomico.
\newblock {\em Rendiconti Lincei}, 3:145--149, 1926.

\bibitem{Pauli1925} Wolfgang Pauli.
\newblock Über den zusammenhang des abschlusses der elektronengruppen im atom
  mit der komplexstruktur der spektren.
\newblock {\em Zeitschrift für Physik}, 31:765--783, 1925.

\bibitem{Pathria2011} R.~K. Pathria and Paul~D. Beale.
\newblock {\em Statistical Mechanics}.
\newblock Elsevier, 3 edition, 2011.

\bibitem{bennett2003} C. H. Bennett, ``Notes on Landauer's principle,
reversible computation, and Maxwell's Demon'', \emph{Studies in History
and Philosophy of Modern Physics} \textbf{34}, 501--510 (2003).
\href{https://doi.org/10.1016/S1355-2198(03)00039-X}{doi:10.1016/S1355-2198(03)00039-X}.
\end{thebibliography}
